\documentclass[12pt]{article}
\usepackage{../../resource/shortex}
\usepackage{mathtools,amssymb}
\usepackage{enumitem}
\usepackage{datetime}
\usepackage{fancyhdr}
\usepackage{ifthen}
\usepackage{pifont}
\usepackage{mathrsfs}
\usepackage[normalem]{ulem}
%\RequirePackage[usenames,dvipsnames]{color}


%% showing solutions 
\newboolean{showsols}
\setboolean{showsols}{false}
\newcommand{\showsolutions}{\setboolean{showsols}{true}}
\newlength{\solspace}
\setlength{\solspace}{16em}
\newcommand{\nosolspace}{\setlength{\solspace}{0em}}

\newcommand{\solution}[2]{\ifthenelse{\boolean{showsols}}{{\par\textcolor{red}{Rubric criterion: #1}\newline\textcolor{blue}{#2}}}{\vspace{\solspace}}}
\newcommand{\npforprint}{\ifthenelse{\boolean{showsols}}{}{\newpage}}
\newcommand{\vsforprint}[1]{\ifthenelse{\boolean{showsols}}{}{\vspace{#1}}}

%% grading
\newcommand{\grades}[1]{~\newline\noindent\framebox[\textwidth][c]{E\quad \quad \quad \quad S \quad \quad \quad \quad R \quad \quad \quad \quad N}}

\newcommand{\questiongrades}{\noindent\emph{Each item is marked $\checkmark$, $\checkmark\!-$, or \ding{55}:}\begin{center}\emph{$\checkmark$ = capable  $~~$ $\checkmark\!-$ = mostly capable  $~~$ \ding{55} = not yet}\end{center} }

\newcommand{\triAssessments}{$\checkmark\quad\checkmark\!-\quad$\ding{55}\quad}
\newcommand{\biAssessments}{$\checkmark\quad$\ding{55}\quad}

%% line for writing name and BUID
\newcommand{\nameBUID}{\noindent\textbf{Name:}\underline{\phantom{XXXXXXXXXXXXXXXXXXX}}\hfill\textbf{BUID:}\underline{\phantom{XXXXXXXXXXXXX}}\newline}


%% formatting 
%\allowdisplaybreaks[3]
\renewcommand{\headrulewidth}{.4mm} % header line width
\renewcommand{\arraystretch}{1.25}


%% semester label, read from the "Semester:" line of important_dates.md
%% (path is relative to a document in materials/ModuleX/<folder>/)
\ExplSyntaxOn
\tl_new:N \g_teaching_semester_tl
\ior_new:N \g_teaching_dates_ior
\file_if_exist:nTF { ../../instructor_inputs/important_dates.md }
  {
    \ior_open:Nn \g_teaching_dates_ior { ../../instructor_inputs/important_dates.md }
    \ior_str_map_inline:Nn \g_teaching_dates_ior
      {
        \regex_extract_once:nnNT { \A \s* Semester: \s* (.*?) \s* \Z } {#1} \l_tmpa_seq
          {
            \tl_gset:Nx \g_teaching_semester_tl { \seq_item:Nn \l_tmpa_seq { 2 } }
            \ior_map_break:
          }
      }
    \ior_close:N \g_teaching_dates_ior
  }
  {
    \msg_new:nnn { teaching } { no-dates }
      { ../../instructor_inputs/important_dates.md~not~found;~semester~header~left~blank. }
    \msg_warning:nn { teaching } { no-dates }
  }
\cs_new:Npn \semester { \tl_use:N \g_teaching_semester_tl }
\ExplSyntaxOff

%% page style
\pagestyle{fancy}
\fancyhf{}
\rhead{\semester}
\lhead{CAS MA 213: Basic Statistics and Probability}
%\rfoot{\thepage}

%--------------------
\begin{document}

\begin{center}
\textbf{\Large Project 2 \\  Statistical Data Analysis }
\end{center}


\section*{Overview}
Your group will select dataset(s) of your interest and conduct one categorical and one numerical data analysis, each including exploratory data analysis (EDA) and a hypothesis test.
\textbf{Your group will write a report of your analyses and submit it together with your R code and data.}

\vspace{6pt}
The project is broken into three deliverables (see the timeline below for the schedule):
\begin{itemize}
    \item \textbf{Deliverable 1 --- Outline:} (P2\_D1\_Outline Google Doc) a written plan covering your project title, group members, background/motivation, data source and description, research questions, variables of interest, categorical and numerical inference plans, data quality, planned visuals, and roles.
    \item \textbf{Deliverable 2 --- Slides and R Code:} (see P2\_D2\_Slides\_and\_Code.pdf for details) draft slides summarizing your results, along with the R code and data used to produce your categorical and numerical analyses.
    \item \textbf{Deliverable 3 --- Written Report:} (see P2\_D3\_Report.pdf for details) your final written report ($\sim$5--10 pages), revised R code, and data, following the suggested outline in \texttt{P2\_D3\_Report.pdf}.
\end{itemize}

\section*{Project requirements}
\begin{itemize}
    \item See the P2\_Rubric.pdf for details and full requirements
    \item Categorical analysis:
    \begin{enumerate}
        \item A description of the research question and the categorical variable(s) of interest
        \item At least one summary table (frequency table or contingency table)
        \item At least one categorical visualization (e.g. bar plot or pie chart)
        \item At least one hypothesis test, including:
        \begin{enumerate}
            \item The null and alternative hypotheses
            \item A check of the conditions for inference
            \item The test statistic and p-value
        \end{enumerate}
        \item A conclusion in the context of the data
    \end{enumerate}
    \item Numerical analysis:
    \begin{enumerate}
        \item A description of the research question and the numerical variable(s) of interest
        \item At least one summary statistics table (mean, median, sd, IQR, etc.)
        \item At least one numerical visualization (e.g. bar plot, histogram, boxplot, or scatterplot)
        \item At least one hypothesis test, including:
        \begin{enumerate}
            \item The null and alternative hypotheses
            \item A check of the conditions for inference
            \item The test statistic and p-value
        \end{enumerate}
        \item A conclusion in the context of the data
    \end{enumerate}
    \item Source of data must be cited, and references should be provided at the end
\end{itemize}

\section*{Hypothesis Tests}
Pick the test that matches your variables and research question.

\subsection*{Categorical Data Analysis (choose at least one)}
\begin{itemize}
    \item \textbf{Single proportion} (Chapter 6.1): one categorical variable with two outcomes. Compare the proportion in one outcome to a hypothesized value.\\
    \emph{Example:} Is the proportion of Titanic passengers who survived different from 50\%?
    \item \textbf{Difference of two proportions} (Chapter 6.2): a two-outcome variable compared between two groups. Test whether the proportion is the same in both groups.\\
    \emph{Example:} Is the survival rate different for female and male Titanic passengers?
    \item \textbf{Chi-square goodness-of-fit} (Chapter 6.3): one categorical variable with three or more outcomes. Test whether the counts match a hypothesized distribution.\\
    \emph{Example:} Are births equally likely on all seven days of the week?
    \item \textbf{Chi-square test of independence} (Chapter 6.4): two categorical variables, at least one with three or more outcomes. Test whether the two variables are associated.\\
    \emph{Example:} Is survival associated with passenger class (1st, 2nd, 3rd) on the Titanic?
\end{itemize}

\subsection*{Numerical Data Analysis (choose at least one)}
\begin{itemize}
    \item \textbf{One-sample mean} (Chapter 7.1): one numerical variable. Compare its mean to a hypothesized value.\\
    \emph{Example:} Do college students sleep less than 8 hours per night on average?
    \item \textbf{Paired means} (Chapter 7.2): two numerical measurements on the same individuals (e.g., before and after). Test whether the mean of the differences is zero.\\
    \emph{Example:} Do students score differently on the final exam than on the midterm, on average?
    \item \textbf{Difference of two means} (Chapter 7.3): one numerical variable compared between two independent groups. Test whether the two group means are equal.\\
    \emph{Example:} Is the average birth weight different for babies of mothers who smoke and mothers who don't?
    \item \textbf{ANOVA} (Chapter 7.5): one numerical variable compared across three or more groups. Test whether all the group means are equal.\\
    \emph{Example:} Does average gas mileage differ for cars with 4, 6, and 8 cylinders?
\end{itemize}

\noindent\textbf{What about regression?} Questions about how one numerical variable changes with another numerical variable are answered with \textbf{linear regression}. Regression is \textbf{not required} for this project (it is not taught until MA214). You may still show a scatterplot of two numerical variables as part of your EDA, but your hypothesis test should be one of the tests above.\\
\emph{Example:} On average, how much does a car's gas mileage change for each additional 1{,}000 pounds of weight? Is this value significantly different from zero? (This is a regression question, not a difference of means question.)

\section*{R Code Guidelines}
\begin{itemize}
    \item You will submit your R code as a separate file (Rscript (.R) or Rmd file), that should be well-commented and organized, along with the data file it reads.
    \item Make sure that your code is reproducible (i.e. that I can run it and get the same results) and that the entire script runs from start to finish without errors. Submit your data file if it is necessary to run your code.
    \item Only use R functions that have been covered in class, tutorials, or labs, or are part of the standard R distribution.
    \item Using certain functions in R is not allowed
    \begin{itemize}
        \item Using functions that directly perform hypothesis tests (e.g., t.test(), chisq.test(), prop.test(), etc.)
        \item If you want to use ANOVA, you can use aov(), anova() or lm() functions.
    \end{itemize}
\end{itemize}

\section*{Writing Guidelines}
\begin{itemize}
    \item Your writing should be clear, concise, and well-structured.
    \item See the GenAI policy on the course website for guidance on using AI tools in your writing.
\end{itemize}

\section*{Possible Data Sources to Consider (not limited to these)}
You can click the links below to explore possible data sources.
\begin{itemize}
    \item \href{https://www.openintro.org/data/}{OpenIntro Datasets}
    \item \href{https://cran.r-project.org/web/packages/fivethirtyeight/vignettes/fivethirtyeight.html}{fivethirtyeight
    R Package}
    \item \href{https://stat.ethz.ch/R-manual/R-devel/library/datasets/html/00Index.html}{datsets R package}
    \item \href{https://www.kaggle.com/datasets}{Kaggle datasets}
    \item \href{https://data.gov/}{US government's open data}
\end{itemize}


\section*{Timeline}
\begin{center}
\begin{tabular}{ll}
\toprule
\textbf{Item} & \textbf{Date} \\
\midrule
Lab P2\_1 (project launch) in class & Fri, Oct 23 \\
\textbf{Deliverable 1 (outline) due} & \textbf{Thu, Oct 29, 9:00 PM} \\
Tutorial 6 due; Lab 6 in class & Fri, Oct 30 \\
Lab P2\_2 (project workday) in class & Fri, Nov 6 \\
\textbf{Deliverable 2 (slides + R code) due} & \textbf{Thu, Nov 12, 9:00 PM} \\
Lab P2\_3 (report prep) in class & Fri, Nov 13 \\
\textbf{Deliverable 3 (report + R code) due} & \textbf{Thu, Nov 19, 9:00 PM} \\
\bottomrule
\end{tabular}
\end{center}



\end{document}
